\documentclass[a4paper,10pt]{article}
\usepackage[utf8]{inputenc}
\usepackage[T1]{fontenc}
\usepackage[french]{babel}
\usepackage{amsmath,amssymb}
\usepackage{geometry}
\usepackage{array,tabularx}
\usepackage{tikz}
\geometry{margin=1.6cm}
\pagestyle{empty}
\newcommand{\ligne}{\par\vspace{0.55cm}\noindent\dotfill\par}
\newcommand{\carreaux}[1]{\begin{center}\begin{tikzpicture}\draw[step=0.5cm,gray!35,very thin] (0,0) grid (17,#1);\end{tikzpicture}\end{center}}
\newenvironment{exercice}[2]{\paragraph{Exercice #1 \hfill #2 points}\mbox{}\par}{\medskip}
\begin{document}
\begin{center}\Large\bfseries Devoir surveillé 17 - Cylindres et volumes\end{center}
\vspace{2mm}
\begin{exercice}{1}{4}Convertir : $3,5$ L en cm$^3$, $4200$ cm$^3$ en L, $0,8$ m$^3$ en L.\end{exercice}\begin{exercice}{2}{5}Calculer le volume d'un pavé droit de dimensions $8$ cm, $5$ cm et $3$ cm.\end{exercice}\begin{exercice}{3}{5}Un cylindre a pour rayon $4$ cm et pour hauteur $10$ cm. Calculer son volume exact puis une valeur approchée au cm$^3$.\end{exercice}\begin{exercice}{4}{6}Un réservoir cylindrique de rayon $30$ cm et de hauteur $80$ cm est rempli aux trois quarts. Quel volume d'eau contient-il en litres ?\end{exercice}\end{document}
